%% Chapter 2 -----------------------------------------------------------------
\chapter{Installation and Startup}
\label{ch:installation}
\index{installation}

Native installers are available for Windows, macOS and Linux. Each bundles the
application with everything it needs to run, including the Python runtime and all
scientific libraries, so no separate Python installation is required. Installers can be
downloaded from the official e-CALLISTO software page,
\url{https://www.e-callisto.org/Software/Callisto-Software.html}, or from the project's
GitHub releases page, \url{https://github.com/SaanDev/e-Callisto_FITS_Analyzer/releases}.

\section{System requirements}
\label{sec:requirements}

\begin{table}[!htbp]
  \centering
  \caption{System requirements for version \appversion.}
  \label{tab:requirements}
  \small
  \begin{tabularx}{\linewidth}{@{}P{0.24\linewidth}L@{}}
    \toprule
    \thead{Item} & \thead{Requirement} \\
    \midrule
    Windows & 64-bit Windows able to run x64 applications. Installation needs administrator approval. \\
    macOS & macOS~15 (Sequoia) or later on a Mac with Apple silicon (arm64). \\
    Linux & Debian, Ubuntu or a derivative on 64-bit x86 (\texttt{amd64}), with the Qt/OpenGL runtime libraries \texttt{libgl1}, \texttt{libegl1}, \texttt{libxkbcommon-x11-0} and \texttt{libxcb-cursor0} (installed automatically by \texttt{apt}). \\
    Disk space & About 3~GB for the application, plus space for the download caches. Solar image sequences can occupy several gigabytes. \\
    Display & 1280~\(\times\)~800 pixels or larger. A resolution of at least 1440~\(\times\)~900 is recommended for the Solar Image Analysis and GCS CME Fitting windows. \\
    Network & Required for the downloaders, archive searches, context viewers, update checks and bug reports. Local files, saved projects and cached downloads can be analyzed offline. \\
    \bottomrule
  \end{tabularx}
\end{table}

\section{Installing on Windows}
\index{installation!Windows}

\begin{steps}
  \item Download the installer, \file{e-CALLISTO_FITS_Analyzer_v3.1.0_Setup.exe}.
  \item Run the installer and approve the Windows administrator prompt.
  \item Follow the installation steps. The installer can optionally create a desktop
    shortcut.
  \item Start the application from the Start menu or the desktop shortcut.
\end{steps}

\begin{note}
  If a save location chosen by the application is not writable --- for example a folder
  inside \file{C:\\Program Files} --- the application asks you to choose another folder.
\end{note}

\section{Installing on macOS}
\index{installation!macOS}

\begin{steps}
  \item Download the disk image (\file{.dmg}) for Apple silicon (arm64).
  \item Open the disk image and drag the application to the \file{Applications} folder.
  \item Eject the disk image and start the application from \file{Applications} or
    Launchpad.
\end{steps}

\begin{note}
  If macOS reports on first launch that the application cannot be opened because its
  developer cannot be verified, open \menu{System Settings > Privacy \& Security} and
  click \gui{Open Anyway}\index{installation!macOS security prompt} for the \appname{}, or Control-click the application in Finder
  and choose \gui{Open}. This is needed only once.
\end{note}

\section{Installing on Linux}
\index{installation!Linux}

Download the Debian package and install it from a terminal in the download folder:
\begin{center}
  \texttt{sudo apt install ./e-callisto-fits-analyzer\_<version>\_amd64.deb}
\end{center}
The package installs a desktop entry, so the application appears in the desktop
environment's application menu. Using \texttt{apt} rather than \texttt{dpkg} also installs
any missing runtime libraries.

\section{Running from source}
\label{sec:from-source}
\index{installation!from source}

The source code is openly available at
\url{https://github.com/SaanDev/e-Callisto_FITS_Analyzer}. Running from source is intended
for developers and for platforms without an installer. With Python~3.12 or newer:
\begin{steps}
  \item Clone or download the repository and create a virtual environment in it.
  \item Install the pinned dependencies with \texttt{python scripts/install\_requirements.py}.
  \item Start the application with \texttt{python src/main.py}.
\end{steps}
On Windows, if Qt reports \texttt{DLL load failed} at start-up, the script
\file{packaging/windows/repair_windows_venv.ps1} rebuilds the virtual environment with the
supported Python version.

\section{Starting the application}
\label{sec:first-launch}

When the application starts it shows a splash screen and then opens the main window
(Chapter~\ref{ch:main-window}) with an empty plotting area. Shortly after start-up it
checks GitHub for a newer release and reports the result on the right of the status bar,
for example \gui{Updates: up to date} (Section~\ref{sec:updates}).

If the previous session ended unexpectedly while it had unsaved work, the application
offers to recover the latest autosave snapshot (Section~\ref{sec:autosave}).

The application remembers your preferences between sessions, including the theme and
rendering mode, which sidebar sections are expanded, the frequency-axis scale, the recent
files and projects, processing and display-range presets, and the JSOC\index{JSOC} e-mail address used
by Solar Image Analysis.

\section{Updating and uninstalling}

Use \menu{About > Check for Updates...}\index{updates!check} to look for a newer version; if one exists, the
installer can be downloaded from within the application (Section~\ref{sec:updates}).
Installing a new version over an existing one keeps your preferences.

To uninstall, use \menu{Settings > Apps} on Windows, move the application from
\file{Applications} to the Trash on macOS, or run
\texttt{sudo apt remove e-callisto-fits-analyzer} on Linux. Downloaded data remain in the
cache folders; empty them first with the \gui{Clear Cache}\index{cache!clearing} buttons of the downloaders if
you want to reclaim the space.
